% Generated draft: requires mathematical and source review.
\documentclass[11pt]{article}
\usepackage[margin=25mm]{geometry}
\usepackage{fontspec}
\setmainfont{Latin Modern Roman}
\usepackage{amsmath,amssymb,mathtools,unicode-math}
\setmathfont{Latin Modern Math}
\usepackage{enumitem,needspace,xurl,hyperref}
\setlength{\parindent}{0pt}
\setlength{\parskip}{5pt}
\setlength{\emergencystretch}{3em}
\newcommand{\N}{\mathbb{N}}
\newcommand{\Z}{\mathbb{Z}}
\newcommand{\Q}{\mathbb{Q}}
\newcommand{\R}{\mathbb{R}}
\newcommand{\C}{\mathbb{C}}
\newcommand{\F}{\mathbb{F}}
\newcommand{\A}{\mathbb{A}}
\newcommand{\PP}{\mathbb{P}}
\newcommand{\E}{\mathbb{E}}
\newcommand{\eps}{\varepsilon}
\newcommand{\dd}{\,\mathrm{d}}
\begin{document}
\section*{UnsolvedMath ID: 700137}
\textbf{Source problem code:} catalogue entry
\subsection{Definitions and mathematical statement}
For a group \(G\), let \(L(G)\) denote the lattice of all subgroups of \(G\), ordered by inclusion, with meet given by intersection and join given by the subgroup generated by the union. For a specified class \(\mathcal H\) of groups, say that \(G\) is determined by its subgroup lattice relative to \(\mathcal H\) if
\[
\forall H\in\mathcal H,\qquad L(H)\cong L(G)\ \Longrightarrow\ H\cong G,
\]
where the first isomorphism is a lattice isomorphism and the second is a group isomorphism. Taking \(\mathcal H\) to be all groups gives the usual unrestricted interpretation; the supplied record does not specify whether a restricted comparison class is intended.

A group \(G\) is nilpotent if its lower central series, defined by \(\gamma_1(G)=G\) and \(\gamma_{n+1}(G)=[\gamma_n(G),G]\), reaches the trivial subgroup after finitely many steps. It is torsion-free if, for every \(g\in G\) and every positive integer \(n\), the equality \(g^n=1\) implies \(g=1\).

The source asks for conditions ensuring that a nilpotent product \(P\) of pure nilpotent groups, drawn from suitable classes, is determined by its subgroup lattice. Thus, after specifying the meaning of pure nilpotent group, the nilpotent-product construction, the permitted factor classes, and the comparison class \(\mathcal H\), the requested conclusion is
\[
\forall H\in\mathcal H,\qquad L(H)\cong L(P)\ \Longrightarrow\ H\cong P.
\]
This is a request to find sufficient conditions, not a single universally quantified conjecture with specified hypotheses. The supplied record does not define purity, specify the number of factors or a nilpotency-class parameter for the product, or identify the intended classes. These omissions remain unresolved.

The source additionally reports that the desired determination is known when the product \(P\) is torsion-free. This is preserved as a supplied background assertion within the source's unspecified setting, rather than asserted here as a verified theorem for every possible interpretation.
\subsection{Short English statement}
Find sufficient conditions under which a nilpotent product of pure nilpotent groups is determined, up to group isomorphism, by its subgroup lattice. The source reports the torsion-free case as known. The meanings of purity, the product construction, the allowed factor classes, and the comparison class require clarification before a fully specified statement can be given.
\subsection{Sources}
\begin{itemize}
\item Local UnsolvedMath registry: permanent ID \(700137\), exact source problem code \texttt{catalogue entry}, title \emph{Kourovka Notebook Problem 2.67}; upstream revision \texttt{d371ee917ac9647fbb06cd6fc529deec6f3dda82}. The registry labels the problem open; this is supplied classification, not an independently verified status assessment.
\item Supplied source attribution: \emph{The Kourovka Notebook (Unsolved Problems in Group Theory)}, original problem 2.67, second issue (1967), as recorded in the twenty-first issue edition (2026), \texttt{arXiv:1401.0300v48}. The collection description names E. I. Khukhro and V. D. Mazurov as editors. No source consultation or independent verification is claimed.
\item Proposed in 1967 by L. E. Sadovski, with an incomplete TeX suffix in both supplied proposer fields. The exact malformed spelling is retained in the accompanying ambiguity record; the intended final letters or diacritics have not been reconstructed.
\item The torsion-free background assertion is supplied without a separate theorem reference, proof, or precise hypotheses. No external URL or separate literature-verification date is supplied.
\end{itemize}
\end{document}
